%%%%%%%%%%%%%%%%%%%%%%%%%%%%%%%%%%%%%%%%%%%%%%%%%%%%%%%%%%%
% --------------------------------------------------------
% Tau
% LaTeX Template
% Version 2.4.4 (28/02/2025)
%
% Author: 
% Guillermo Jimenez (memo.notess1@gmail.com)
% 
% License:
% Creative Commons CC BY 4.0
% --------------------------------------------------------
%%%%%%%%%%%%%%%%%%%%%%%%%%%%%%%%%%%%%%%%%%%%%%%%%%%%%%%%%%%

\documentclass[9pt,a4paper,twocolumn,twoside]{tau-class/tau}
\usepackage[english]{babel}

%% Spanish babel recomendation
% \usepackage[spanish,es-nodecimaldot,es-noindentfirst]{babel} 

%% Draft watermark
% \usepackage{draftwatermark}

%----------------------------------------------------------
% TITLE
%----------------------------------------------------------

\journalname{Rapporto tecnico}
\title{Rilievo fotogrammetrico della spiaggia di Borghetto S.S. (SV)}

%----------------------------------------------------------
% AUTHORS, AFFILIATIONS AND PROFESSOR
%----------------------------------------------------------

\author[a]{Alessio Rovere}

%----------------------------------------------------------

\professor{GeoDrone s.r.l, Via Torino 155 - 30172 Venezia}

%----------------------------------------------------------
% FOOTER INFORMATION
%----------------------------------------------------------

\footinfo{Esercitazione corso Geografia Fisica e Geomorfologia}
\course{Creative Commons CC BY 4.0}

%----------------------------------------------------------
% ABSTRACT AND KEYWORDS
%----------------------------------------------------------
\begin{abstract}    
Su incarico del gruppo di ricerca del corso di \textit{Geografia Fisica e Geomorfologia}, la ditta \textbf{GeoDrone S.r.l.} ha eseguito nel corso del 2024 un rilievo fotogrammetrico dell’area costiera di \textbf{Borghetto Santo Spirito (SV)}, con l’obiettivo di ottenere un \textit{Modello Digitale di Elevazione (DEM)} e un’\textit{ortofoto ad alta risoluzione} da utilizzare per analisi geomorfologiche e studi di impatto da mareggiata.
\end{abstract}

%----------------------------------------------------------

\begin{document}
		
    \maketitle 
    \thispagestyle{firststyle} 
    \tauabstract 
    % \tableofcontents
    % \linenumbers 
    
%----------------------------------------------------------

\section*{Metodologia di rilievo}

La fotogrammetria da drone è una tecnica di rilievo che permette di ricostruire in 3D superfici e morfologie del terreno a partire da immagini digitali acquisite da piattaforma aerea. 
Il principio di base è la \textbf{visione stereoscopica}: un medesimo punto viene ripreso da più angolazioni, e la differenza di posizione apparente (parallasse) tra le immagini consente di calcolarne la posizione tridimensionale mediante triangolazione. 

Rispetto ai rilievi topografici tradizionali, questa metodologia consente di ottenere una copertura molto più densa del terreno, con tempi di acquisizione ridotti e risoluzioni planimetriche e altimetriche dell’ordine dei centimetri. 

Il rilievo di Borghetto è stato condotto mediante drone multirotore \textit{DJI Mavic Mini 2} equipaggiato con fotocamera \textbf{DJI FC7303 (4.49 mm)}.  
Sono state acquisite complessivamente \textbf{1.502 immagini} ad un’altezza media di volo di \textbf{47.1 m}, con una \textbf{Ground Sampling Distance (GSD)} di \textbf{1.45 cm/pixel} e una copertura totale di circa \textbf{0.104 km$^2$}.  

La pianificazione del volo ha previsto un’elevata sovrapposizione longitudinale e trasversale delle immagini (circa 80\%), al fine di garantire una ridondanza ottimale dei punti omologhi e migliorare la qualità della ricostruzione tridimensionale.  
La Figura~1 del report di elaborazione mostra la distribuzione delle immagini e l’area coperta dal rilievo.

%----------------------------------------------------------

\section*{Punti di controllo e accuratezza}

Per ottenere una georeferenziazione accurata dei prodotti fotogrammetrici, sono stati utilizzati \textbf{18 Ground Control Points (GCP)}, rilevati con strumentazione GNSS \textbf{EMLID REACH RS2} operante in configurazione \textit{base–rover}, in coordinate \textbf{WGS84 / UTM zone 32N (EPSG:32632)}. Le altezze sono state riferite al geoide ITG 2009 \cite{corchete2010high}.

Nel sistema \textbf{base–rover}, un ricevitore (base) rimane fisso su un punto noto e trasmette correzioni differenziali in tempo reale a un secondo ricevitore (rover) che effettua le misure sui punti da determinare. 
Questa tecnica, detta \textbf{posizionamento differenziale GNSS}, riduce in modo significativo gli errori dovuti all’atmosfera e alle orbite satellitari, consentendo precisioni planimetriche e altimetriche dell’ordine del centimetro.

Le coordinate dei GCP sono state importate in Metashape e utilizzate per la calibrazione e l’ottimizzazione del modello, garantendo una scala e un posizionamento accurati.  

L’errore medio quadratico totale (RMSE) risulta pari a \textbf{4.7 cm}, con valori medi di:
\begin{itemize}
    \item X error: 2.7 cm
    \item Y error: 2.2 cm
    \item Z error: 3.1 cm
\end{itemize}

Questa accuratezza è pienamente conforme agli standard per rilievi fotogrammetrici di dettaglio a scala locale.

%----------------------------------------------------------

\section*{Principi di elaborazione fotogrammetrica}

Il processo di elaborazione delle immagini è stato eseguito nel software \textbf{Agisoft Metashape Professional v2.2.2}, seguendo un workflow standardizzato articolato nelle seguenti fasi:

\begin{enumerate}
    \item \textbf{Allineamento delle immagini:} identificazione automatica di punti omologhi (\textit{tie points}) e stima della posizione e dell’orientamento delle camere.
    \item \textbf{Ottimizzazione e calibrazione:} correzione delle distorsioni ottiche e affinamento dei parametri interni della fotocamera (focale, centro ottico, distorsioni radiali e tangenziali).
    \item \textbf{Generazione della nuvola di punti densa:} ricostruzione tridimensionale dell’area mediante confronto pixel–per–pixel tra immagini corrispondenti (\textit{depth maps}).
    \item \textbf{Creazione del modello 3D (mesh):} interpolazione dei punti per formare una superficie continua.
    \item \textbf{Produzione di DEM e ortofoto:} proiezione ortogonale dei dati e generazione di prodotti georeferenziati.
\end{enumerate}

La fotogrammetria digitale si basa su tecniche di triangolazione e sulla geometria delle immagini.  
A partire da due o più fotografie che ritraggono lo stesso punto del terreno, è possibile calcolare la posizione tridimensionale del punto nello spazio.  
L’insieme di questi punti costituisce la \textbf{nuvola di punti densa}, dalla quale vengono derivati i prodotti finali: DEM e ortofoto.

%----------------------------------------------------------

\section*{Risultati}

Il processo di ricostruzione ha generato i seguenti prodotti principali:
\begin{itemize}
    \item \textbf{Modello Digitale di Elevazione (DEM)} con risoluzione di \textbf{2.90 cm/pixel} e densità di \textbf{0.119 punti/cm$^2$};
    \item \textbf{Ortofoto} con risoluzione di \textbf{10 cm/pixel}, ottenuta tramite mosaicatura e filtraggio per la rimozione di artefatti.
\end{itemize}

Entrambi i prodotti sono stati successivamente ridimensionati per ottenere una risoluzione omogenea di \textbf{10 cm/pixel}, così da facilitarne l’integrazione nei software GIS.  

Tutti i modelli sono stati referenziati nel sistema \textbf{WGS84 / UTM 32N}, per garantire compatibilità con ambienti di analisi spaziale come \textit{QGIS} e \textit{ArcGIS}.

%----------------------------------------------------------

\section*{Considerazioni finali}

L’insieme dei prodotti ottenuti (ortofoto e DEM) costituisce una base metrica solida per analisi di dettaglio sull’evoluzione della linea di riva, il calcolo di profili topografici e l’individuazione di aree soggette a inondazione o erosione.  

Il rilievo fotogrammetrico, integrato con tecniche GNSS ad alta precisione, rappresenta oggi uno strumento fondamentale per lo studio delle dinamiche costiere e la pianificazione degli interventi di gestione e mitigazione del rischio.  
Le studentesse e gli studenti possono utilizzare i dati forniti in ambiente GIS per sovrapporre modelli di runup e identificare le zone di massima inondazione, come previsto nelle esercitazioni successive.

%----------------------------------------------------------

\printbibliography

%----------------------------------------------------------

\clearpage

\begin{tauenv}[frametitle=Citazione del report tecnico e dei dati associati]
Per citare correttamente il report del rilievo e i dati forniti da \textbf{GeoDrone S.r.l.}, è possibile aggiungere la seguente voce BibTeX in \textbf{Zotero}.  
Zotero permette di gestire automaticamente citazioni e bibliografia anche in \textbf{Microsoft Word}, utilizzando lo stile scelto (APA, Harvard, ecc.).

\vspace{0.5em}
\noindent\begin{minipage}{\columnwidth}
\scriptsize
\raggedright
\begin{verbatim}
@techreport{GeoDrone2025_Borghetto,
  author       = {Rovere, Alessio and Casella, Emanuele 
                  and Meneghetti, Nicola},
  title        = {Rilievo fotogrammetrico della spiaggia 
                  di Borghetto S.S. (SV)},
  institution  = {GeoDrone S.r.l.},
  address      = {Via Torino 155, 30172 Venezia, Italia},
  year         = {2025},
  month        = {October},
  note         = {Rapporto tecnico. Elaborazione dei dati 
                  fotogrammetrici con Agisoft Metashape 
                  Professional v2.2.2.},
  url          = {https://doi.org/10.12345/
                  geodrone.borghetto.2025},
  type         = {Technical Report}
}
\end{verbatim}
\end{minipage}

\vspace{0.5em}
\noindent
\textbf{Come importare la citazione in Zotero:}
\begin{enumerate}
  \item Copiare tutto il testo del blocco BibTeX qui sopra;
  \item In Zotero, selezionare: \texttt{File → Import → From Clipboard};
  \item Zotero riconoscerà automaticamente la voce come \textit{Technical Report};
  \item Controllare che i campi (autori, titolo, anno, istituzione) siano corretti e, se necessario, aggiungere tag come “Fotogrammetria” o “Borghetto” per organizzare la libreria.
\end{enumerate}

\noindent
\textbf{Come usare la citazione in Word:}
\begin{enumerate}
  \item Aprire il documento in Word;
  \item Nel menu Zotero (scheda superiore), posizionare il cursore dove si vuole inserire la citazione e cliccare su \texttt{Add/Edit Citation};
  \item Cercare “GeoDrone 2025 Borghetto” e selezionare la voce;
  \item A fine documento, cliccare su \texttt{Add/Edit Bibliography} per inserire automaticamente la bibliografia.
\end{enumerate}

Zotero formatterà automaticamente citazioni e bibliografia in base allo stile scelto (APA, Chicago, Harvard, ecc.).
\end{tauenv}





\end{document}